% !TEX program = xelatex
% ============================================================================
%  POSN Computer Qualification — Sequences Slide Deck
%  Topic: ลำดับ (Sequences)
% ============================================================================

\documentclass[11pt, aspectratio=169]{beamer}

% --- Load shared preamble ---
% ============================================================================
%  POSN Computer Qualification — Shared Beamer Preamble
%  Used by all topic slide decks. Compile with XeLaTeX.
% ============================================================================

% ---------- Thai language & font ----------
\usepackage{iftex}
\usepackage{fontspec}

% XeTeX-specific: enable Thai line breaking
\ifXeTeX
  \XeTeXlinebreaklocale{th}
\fi
% Noto Sans Thai — body text (clean modern sans-serif, handles all Thai)
\setmainfont{Noto Sans Thai}[
  UprightFont = Noto Sans Thai Light,
  BoldFont    = Noto Sans Thai Medium,
  Scale       = 0.9
]
\setsansfont{Noto Sans Thai}[
  UprightFont = Noto Sans Thai Light,
  BoldFont    = Noto Sans Thai Medium,
  Scale       = 0.9
]

% Noto Sans Thai — clean modern sans-serif for structural elements
\newfontfamily\headingfont{Noto Sans Thai}[
  UprightFont = Noto Sans Thai Medium,
  BoldFont    = Noto Sans Thai Bold,
  Scale       = 0.9
]

% --- Heading font for structural elements ---
\setbeamerfont{title}{family=\headingfont, series=\bfseries, size=\huge}
\setbeamerfont{subtitle}{family=\headingfont, series=\mdseries}
\setbeamerfont{author}{family=\headingfont, series=\mdseries}

% --- Custom Title Page ---
\defbeamertemplate{title page}{custom}[1][]{%
  \centering
  \vspace*{2cm}
  % Title — in white
  {\usebeamerfont{title}\color{white}\inserttitle\par}
  \vspace{0.6cm}
  % Subtitle — in white
  {\usebeamerfont{subtitle}\color{white}\insertsubtitle\par}
  \vspace{2.5cm}
  % Author — blue filled rounded box, white text
  \begin{center}
    \begin{tcolorbox}[arc=3mm, colback=boxBlue, colframe=boxBlue!80,
                      width=0.42\textwidth, halign=center,
                      left=1.5em, right=1.5em, top=0.8em, bottom=0.8em,
                      boxrule=1.5pt, coltext=white]
      \usebeamerfont{author}\insertauthor
    \end{tcolorbox}
  \end{center}
}
\setbeamertemplate{title page}[custom]
\setbeamerfont{frametitle}{family=\headingfont, series=\bfseries}
\setbeamerfont{section in toc}{family=\headingfont}
\setbeamerfont{page number in head/foot}{family=\headingfont}
\setbeamerfont{section title}{family=\headingfont, series=\bfseries}
\setbeamerfont{subsection title}{family=\headingfont}

% Monospace font (for code/verbatim)
\setmonofont{Latin Modern Mono}[Scale=0.9]

% ---------- Math ----------
\usepackage{amsmath, amssymb}

% ---------- Graphics ----------
\usepackage{graphicx}
\usepackage{tikz}

% ---------- String parsing (for section headers) ----------
\usepackage{xstring}

% ---------- Colored boxes ----------
\usepackage[skins, breakable, xparse]{tcolorbox}
% Note: we do NOT load the 'most' option — it predefines example/solution environments
% that would conflict with our custom boxes.

% --- Color definitions ---
\definecolor{boxBlue}{RGB}{41, 128, 185}
\definecolor{boxGreen}{RGB}{39, 174, 96}
\definecolor{boxOrange}{RGB}{230, 126, 34}
\definecolor{boxPurple}{RGB}{142, 68, 173}
\definecolor{boxBlueBg}{RGB}{235, 245, 255}
\definecolor{boxGreenBg}{RGB}{235, 255, 240}
\definecolor{boxOrangeBg}{RGB}{255, 245, 235}
\definecolor{boxPurpleBg}{RGB}{248, 240, 255}

% --- Explanation box (blue) — for definitions, concepts, notes ---
\newtcolorbox{explanation}[1][]{
  enhanced,
  breakable,
  arc         = 2mm,
  colback     = boxBlueBg,
  colframe    = boxBlue!50,
  title       = #1,
  fonttitle   = \bfseries\color{white},
  coltitle    = white,
  colbacktitle= boxBlue,
  attach boxed title to top left = {yshift=-2mm, xshift=2mm},
  boxed title style = {arc=2mm, size=small},
  before skip = 8pt,
  after skip  = 8pt
}

% --- Example box (green) — for worked examples ---
\newtcolorbox{exambox}[1][]{
  enhanced,
  breakable,
  arc         = 2mm,
  colback     = boxGreenBg,
  colframe    = boxGreen!50,
  title       = #1,
  fonttitle   = \bfseries\color{white},
  coltitle    = white,
  colbacktitle= boxGreen,
  attach boxed title to top left = {yshift=-2mm, xshift=2mm},
  boxed title style = {arc=2mm, size=small},
  before skip = 8pt,
  after skip  = 8pt
}

% --- Solution box (orange) — for solutions and answers ---
\newtcolorbox{solbox}[1][]{
  enhanced,
  breakable,
  arc         = 2mm,
  colback     = boxOrangeBg,
  colframe    = boxOrange!50,
  title       = #1,
  fonttitle   = \bfseries\color{white},
  coltitle    = white,
  colbacktitle= boxOrange,
  attach boxed title to top left = {yshift=-2mm, xshift=2mm},
  boxed title style = {arc=2mm, size=small},
  before skip = 8pt,
  after skip  = 8pt
}

% --- Intuition box (purple) — for plain-language explanations, analogies ---
\newtcolorbox{intuition}[1][]{
  enhanced,
  breakable,
  arc         = 2mm,
  colback     = boxPurpleBg,
  colframe    = boxPurple!50,
  title       = #1,
  fonttitle   = \bfseries\color{white},
  coltitle    = white,
  colbacktitle= boxPurple,
  attach boxed title to top left = {yshift=-2mm, xshift=2mm},
  boxed title style = {arc=2mm, size=small},
  before skip = 8pt,
  after skip  = 8pt
}

% ---------- Beamer theme ----------
\usetheme{Madrid}
\usecolortheme{default}

% --- Custom beamer colors (blue accent matching explanation boxes) ---
\setbeamercolor{structure}{fg=boxBlue}
\setbeamercolor{palette primary}{fg=white, bg=boxBlue}
\setbeamercolor{palette secondary}{fg=white, bg=boxBlue!80}
\setbeamercolor{palette tertiary}{fg=white, bg=boxBlue!60}
\setbeamercolor{block title}{fg=white, bg=boxBlue}
\setbeamercolor{block body}{fg=black, bg=boxBlueBg}

% Remove navigation symbols for clean look
\setbeamertemplate{navigation symbols}{}

% Note: Frame title prefix removed — \addtobeamertemplate conflicts with beamer internals

% --- Outline (Table of Contents) styling ---
\setbeamerfont{section in toc}{family=\headingfont, series=\mdseries}
\setbeamerfont{subsection in toc}{family=\headingfont, series=\mdseries}
\setbeamertemplate{section in toc}{%
  \leavevmode\leftskip=1.5em
  \textcolor{boxBlue}{\textbf{\large\inserttocsectionnumber.}}%
  \quad\inserttocsection\par
  \vspace{0.2em}
}
\setbeamertemplate{subsection in toc}{%
  \leavevmode\leftskip=4em
  \textcolor{boxBlue!60}{\small$\bullet$}%
  \quad\inserttocsubsection\par
  \vspace{0.1em}
}

% Section header slides — Thai in white, English in smaller light blue below
\AtBeginSection{
  \begin{frame}[plain]{}
    \vspace*{2cm}
    \begin{beamercolorbox}[sep=12pt, center, rounded=true, shadow=true]{palette primary}
      \IfSubStr{\secname}{(}{%
        \StrBefore{\secname}{(}[\sectionthai]%
        \StrBetween[1,1]{\secname}{(}{)}[\sectioneng]%
        \begin{tabular}{@{}c@{}}
          {\usebeamerfont{title}\sectionthai} \\[0.2em]
          {\usebeamerfont{subtitle}\color{boxBlue!60}\sectioneng}
        \end{tabular}%
      }{%
        {\usebeamerfont{title}\secname\par}%
      }
    \end{beamercolorbox}
    \vfill
  \end{frame}
}

% Subsection header slides — smaller, centered
\AtBeginSubsection{
  \begin{frame}[plain]{}
    \vfill
    \centering
    {\usebeamerfont{section title}\insertsubsectionhead\par}
    \vfill
  \end{frame}
}

% Minimal footer: page number only, right-aligned
\setbeamertemplate{footline}{
  \hfill\usebeamerfont{page number in head/foot}\insertframenumber{} / \inserttotalframenumber\hspace*{1em}\vspace*{0.5ex}
}

% ---------- Spacing & readability ----------
\linespread{1.2}

% ---------- Useful commands ----------

% Highlight important text in blue
\newcommand{\highlight}[1]{\textcolor{boxBlue}{\textbf{#1}}}

% Math set notation helpers
\newcommand{\R}{\mathbb{R}}
\newcommand{\N}{\mathbb{N}}
\newcommand{\Z}{\mathbb{Z}}
\newcommand{\Q}{\mathbb{Q}}

% ============================================================================
%  End of preamble
% ============================================================================


% --- Document info ---
\title[ลำดับ]{ลำดับ}
\subtitle{Sequences — ลำดับเลขคณิต, ลำดับเรขาคณิต, ลำดับฟิโบนักชี}
\author[None W. (Yuan)]{None W. (Yuan)}
\date{}

\begin{document}

% ===========================================================================
%  TITLE SLIDE
% ===========================================================================
\begin{frame}[plain]
  \titlepage
\end{frame}

% ===========================================================================
%  OUTLINE
% ===========================================================================
\begin{frame}{สารบัญ (Outline)}
  \tableofcontents
\end{frame}

% ===========================================================================
%  ก่อนเรียน (Prerequisites)
% ===========================================================================
\begin{frame}{ก่อนเรียน (Prerequisites)}
  \begin{explanation}[สิ่งที่ควรรู้ก่อนเรียน]
    \begin{itemize}
      \item \textbf{จำนวนจริง} — การดำเนินการพื้นฐาน, เลขยกกำลัง
      \item \textbf{ฟังก์ชัน} — แนวคิดของฟังก์ชันที่มีโดเมนเป็นจำนวนนับ
    \end{itemize}
  \end{explanation}

  \begin{intuition}[ลำดับคืออะไร?]
    \textbf{ลำดับ} คือฟังก์ชันที่มีโดเมนเป็นจำนวนเต็มบวก $\{1, 2, 3, \ldots\}$
    เรียก $a_n$ ว่า \highlight{พจน์ที่ $n$} ของลำดับ

    \medskip
    เขียนแทนด้วย $a_1, a_2, a_3, \ldots$ หรือ $\{a_n\}$
  \end{intuition}
\end{frame}

% ===========================================================================
%  SECTION 1: ลำดับเลขคณิต (Arithmetic Sequence)
% ===========================================================================
\section{ลำดับเลขคณิต (Arithmetic Sequence)}

\subsection{นิยามและสูตร}

\begin{frame}{ลำดับเลขคณิตคืออะไร?}
  \begin{explanation}[นิยาม]
    \textbf{ลำดับเลขคณิต} คือลำดับที่ \highlight{ผลต่างระหว่างพจน์ที่อยู่ติดกันมีค่าคงที่}
    เรียกค่าคงที่นี้ว่า \textbf{ผลต่างร่วม} (common difference) เขียนแทนด้วย $d$

    \medskip
    $d = a_{n+1} - a_n$ \quad (ค่าคงที่ สำหรับทุก $n$)
  \end{explanation}

  \begin{exambox}[ตัวอย่าง]
    $2, 5, 8, 11, 14, \ldots$

    \medskip
    $d = 5 - 2 = 3, \quad 8 - 5 = 3, \quad 11 - 8 = 3, \ldots$

    \medskip
    เป็นลำดับเลขคณิตที่มี $a_1 = 2$, $d = 3$
  \end{exambox}
\end{frame}

\begin{frame}{เลขคณิต vs เรขาคณิต: จำง่าย ๆ}
  \begin{intuition}[จำง่าย ๆ]
    \begin{itemize}
      \item \textbf{เลขคณิต} = \highlight{บวกเพิ่มทีละเท่า ๆ กัน} (หรือลบออก)

      \quad เช่น $3, 7, 11, 15, \ldots$ (บวก 4 ตลอด)

      \item \textbf{เรขาคณิต} = \highlight{คูณเพิ่มทีละเท่า ๆ กัน} (หรือหาร)

      \quad เช่น $3, 6, 12, 24, \ldots$ (คูณ 2 ตลอด)
    \end{itemize}
  \end{intuition}
\end{frame}

\begin{frame}{พจน์ทั่วไปของลำดับเลขคณิต}
  \begin{explanation}[สูตรพจน์ที่ $n$]
    \[
      a_n = a_1 + (n-1)d
    \]

    \textbf{ความหมาย:} เริ่มจาก $a_1$ แล้วบวก $d$ เพิ่มอีก $(n-1)$ ครั้ง
  \end{explanation}

  \begin{exambox}[ตัวอย่าง: หาพจน์ที่ 20]
    จงหาพจน์ที่ 20 ของลำดับ $3, 7, 11, 15, \ldots$

    \medskip
    $a_1 = 3$, $d = 4$

    $a_{20} = 3 + (20-1) \times 4 = 3 + 76 = 79$
  \end{exambox}
\end{frame}

% \begin{frame}{ผลบวกของพจน์ในลำดับเลขคณิต}
%   \begin{explanation}[สูตรผลบวก $n$ พจน์แรก]
%     \[
%       S_n = \frac{n}{2}\bigl(a_1 + a_n\bigr)
%     \]

%     หรือเขียนในรูปของ $a_1$ และ $d$:
%     \[
%       S_n = \frac{n}{2}\bigl[2a_1 + (n-1)d\bigr]
%     \]
%   \end{explanation}

%   \begin{exambox}[ตัวอย่าง]
%     จงหาผลบวกของ $1 + 2 + 3 + \cdots + 100$

%     \medskip
%     $a_1 = 1$, $a_{100} = 100$, $n = 100$

%     $S_{100} = \frac{100}{2}(1 + 100) = 50 \times 101 = 5050$
%   \end{exambox}
% \end{frame}

% ===========================================================================
%  SECTION 2: ลำดับเรขาคณิต (Geometric Sequence)
% ===========================================================================
\section{ลำดับเรขาคณิต (Geometric Sequence)}

\subsection{นิยามและสูตร}

\begin{frame}{ลำดับเรขาคณิตคืออะไร?}
  \begin{explanation}[นิยาม]
    \textbf{ลำดับเรขาคณิต} คือลำดับที่ \highlight{อัตราส่วนระหว่างพจน์ที่อยู่ติดกันมีค่าคงที่}
    เรียกค่าคงที่นี้ว่า \textbf{อัตราส่วนร่วม} (common ratio) เขียนแทนด้วย $r$

    \medskip
    $r = \dfrac{a_{n+1}}{a_n}$ \quad (ค่าคงที่ สำหรับทุก $n$)
  \end{explanation}

  \begin{exambox}[ตัวอย่าง]
    $3, 6, 12, 24, 48, \ldots$

    \medskip
    $r = \dfrac{6}{3} = 2, \quad \dfrac{12}{6} = 2, \quad \dfrac{24}{12} = 2, \ldots$

    \medskip
    เป็นลำดับเรขาคณิตที่มี $a_1 = 3$, $r = 2$
  \end{exambox}
\end{frame}

\begin{frame}{พจน์ทั่วไปของลำดับเรขาคณิต}
  \begin{explanation}[สูตรพจน์ที่ $n$]
    \[
      a_n = a_1 \cdot r^{\,n-1}
    \]

    \textbf{ความหมาย:} เริ่มจาก $a_1$ แล้วคูณด้วย $r$ อีก $(n-1)$ ครั้ง
  \end{explanation}

  \begin{exambox}[ตัวอย่าง: หาพจน์ที่ 6]
    จงหาพจน์ที่ 6 ของลำดับ $2, 6, 18, 54, \ldots$

    \medskip
    $a_1 = 2$, $r = 3$

    $a_6 = 2 \times 3^{5} = 2 \times 243 = 486$
  \end{exambox}
\end{frame}

% \begin{frame}{ผลบวกของพจน์ในลำดับเรขาคณิต}
%   \begin{explanation}[สูตรผลบวก $n$ พจน์แรก]
%     \[
%       S_n = \frac{a_1(1 - r^n)}{1 - r} \qquad (r \neq 1)
%     \]

%     หรือเขียนอีกแบบ: \quad
%     $S_n = \dfrac{a_1(r^n - 1)}{r - 1}$
%   \end{explanation}

%   \begin{exambox}[ตัวอย่าง]
%     จงหาผลบวก 5 พจน์แรกของลำดับ $3, 6, 12, 24, \ldots$

%     \medskip
%     $a_1 = 3$, $r = 2$, $n = 5$

%     $S_5 = \dfrac{3(1 - 2^5)}{1 - 2}
%          = \dfrac{3(1 - 32)}{-1}
%          = \dfrac{3 \times (-31)}{-1}
%          = 93$
%   \end{exambox}
% \end{frame}

\begin{frame}{เปรียบเทียบ: เลขคณิต vs เรขาคณิต}
  \begin{columns}[T]
    \column{0.48\textwidth}
      \textbf{ลำดับเลขคณิต (Arithmetic)}
      \begin{itemize}
        \item \highlight{บวก} เพิ่มทีละ $d$
        \item $a_n = a_1 + (n-1)d$
        % \item $S_n = \frac{n}{2}(a_1 + a_n)$
        \item เช่น $2, 5, 8, 11, \ldots$
      \end{itemize}

    \column{0.48\textwidth}
      \textbf{ลำดับเรขาคณิต (Geometric)}
      \begin{itemize}
        \item \highlight{คูณ} เพิ่มทีละ $r$
        \item $a_n = a_1 \cdot r^{\,n-1}$
        % \item $S_n = \frac{a_1(1-r^n)}{1-r}$
        \item เช่น $2, 6, 18, 54, \ldots$
      \end{itemize}
  \end{columns}

  \begin{intuition}[จุดแตกต่างสำคัญ]
    เลขคณิต = เพิ่มแบบ \textbf{เส้นตรง} (linear) \quad
    เรขาคณิต = เพิ่มแบบ \textbf{เอกซ์โพเนนเชียล} (exponential)
  \end{intuition}
\end{frame}

% ===========================================================================
%  SECTION 3: ลำดับฟิโบนักชี (Fibonacci Sequence)
% ===========================================================================
\section{ลำดับฟิโบนักชี (Fibonacci Sequence)}

\begin{frame}{ลำดับฟิโบนักชีคืออะไร?}
  \begin{explanation}[นิยาม]
    \textbf{ลำดับฟิโบนักชี} คือลำดับที่แต่ละพจน์ (ตั้งแต่พจน์ที่ 3 เป็นต้นไป)
    เท่ากับ \highlight{ผลบวกของสองพจน์ก่อนหน้า}

    \medskip
    $F_1 = 1, \quad F_2 = 1$

    $F_n = F_{n-1} + F_{n-2}$ \quad สำหรับ $n \geq 3$
  \end{explanation}

  \begin{exambox}[10 พจน์แรกของลำดับฟิโบนักชี]
    \[
      1, 1, 2, 3, 5, 8, 13, 21, 34, 55, \ldots
    \]

    \medskip
    \begin{itemize}
      \item $F_3 = F_2 + F_1 = 1 + 1 = 2$
      \item $F_4 = F_3 + F_2 = 2 + 1 = 3$
      \item $F_5 = F_4 + F_3 = 3 + 2 = 5$
    \end{itemize}
  \end{exambox}
\end{frame}

% ===========================================================================
%  SUMMARY
% ===========================================================================
\section{สรุป}

\begin{frame}{สรุปเนื้อหา}
  \begin{explanation}[สิ่งที่ได้เรียนรู้]
    \begin{enumerate}
      \item \textbf{ลำดับเลขคณิต}: ผลต่างร่วม $d$ คงที่
        \quad $a_n = a_1 + (n-1)d$
        % \quad $S_n = \frac{n}{2}(a_1 + a_n)$

      \item \textbf{ลำดับเรขาคณิต}: อัตราส่วนร่วม $r$ คงที่
        \quad $a_n = a_1 \cdot r^{\,n-1}$
        % \quad $S_n = \frac{a_1(1-r^n)}{1-r}$

      \item \textbf{ลำดับฟิโบนักชี}: $F_1 = 1, F_2 = 1$,
        $F_n = F_{n-1} + F_{n-2}$
    \end{enumerate}
  \end{explanation}

  \begin{center}
    \large เอกสารนี้เป็นส่วนหนึ่งของเนื้อหาเตรียมสอบ \textbf{สอวน. คอมพิวเตอร์}
  \end{center}
\end{frame}

\end{document}
